\documentclass[preprint,12pt]{elsarticle}

\usepackage[utf8]{inputenc}
\usepackage[english]{babel}
\usepackage{graphicx}
\usepackage{booktabs}
\usepackage{amsmath}
\usepackage{amsfonts}
\usepackage{amssymb}
\usepackage{hyperref}
\usepackage{xurl}
\usepackage{ragged2e}
\usepackage{array}
\usepackage{xcolor}
\usepackage{geometry}
\usepackage{longtable}
\usepackage{microtype}
\usepackage{todonotes}

\newcommand{\mycomment}[1]{\textcolor{red}{[JGF: #1]}}
\geometry{a4paper, margin=2.5cm}

\newcolumntype{P}[1]{>{\RaggedRight\arraybackslash}p{#1}}

\journal{Journal of Informetrics}

\begin{document}

\begin{frontmatter}

\title{OpenSciencePipe: An End-to-End Computational Framework for Autonomous Science Mapping}

\author[inst1]{João Garcia Farinha\corref{cor1}}
\ead{jgfarinha@ciencias.ulisboa.pt}
\cortext[cor1]{Corresponding author}
\author[inst1]{Carlos Duarte}
\ead{cduarte@ciencias.ulisboa.pt}

\affiliation[inst1]{organization={LASIGE, Faculdade de Ciências, Universidade de Lisboa},
            addressline={Campo Grande}, 
            city={Lisboa},
            postcode={1749-016},
            country={Portugal}}

\begin{abstract}
Traditional bibliometric science mapping workflows frequently rely on desktop applications such as VOSviewer, CiteSpace, and Bibliometrix, which require manual data export from commercial databases, depend on keyword co-occurrence heuristics that discard semantic context, and struggle to scale computationally to large literature corpora. In this paper, we present \textbf{OpenSciencePipe}, an open-source, computational framework for autonomous bibliometric discovery and science mapping. The framework integrates multi-database ingestion across OpenAlex, Crossref, Semantic Scholar, and PubMed, programmatic deduplication, GPU-accelerated co-authorship and coupling network topologies, unsupervised neural thematic extraction via BERTopic sentence embeddings with HDBSCAN density-based noise rejection, and local Large Language Model zero-shot classification via Ollama for deductive paradigm extraction and systematic relevance screening. We evaluate OpenSciencePipe against a ground-truth suite of 20 published bibliometric benchmark studies spanning medicine, artificial intelligence, econometrics, environmental sciences, and materials engineering. The neural clustering architecture achieves high topic diversity (mean inverted Rank-Biased Overlap $\bar{\text{RBO}} = 0.79$, spanning $0.62 \le \text{RBO} \le 0.95$) and prevents semantic boundary contamination by isolating peripheral literature into an explicit Topic~$-1$ noise pool. Furthermore, the automated pipeline exhibits strong topological concordance with published bibliometric structures, replicating geopolitical country collaboration networks with near-perfect ordinal alignment, achieving up to $100.0\%$ Top-15 country overlap with Spearman correlation $\rho \ge 0.96$ ($p < 0.001$). GPU acceleration delivers sub-second graph community partitioning on large-scale bibliometric graphs with $N > 50,000$. By automating the entire trajectory from API query execution to publication-ready synthesis, OpenSciencePipe establishes a scalable, reproducible foundation for next-generation informetric research.
\end{abstract} 


\begin{highlights}
\item OpenSciencePipe: An end-to-end autonomous science mapping pipeline unifying API ingestion, GPU graphs, and hybrid NLP.
\item Empirical validation against 20 published benchmark studies across multiple scientific disciplines.
\item Unsupervised BERTopic modeling achieves high topic diversity ($\bar{\text{RBO}} = 0.79$) with HDBSCAN Topic~$-1$ noise rejection.
\item Local Ollama LLM integration enables zero-shot relevance screening and deductive paradigm extraction without cloud API exposure.
\item Sub-second GPU-accelerated Louvain community detection and PageRank scaling on corpora exceeding 50,000 papers.
\item Significant topological concordance with established human and VOSviewer benchmarks ($\rho \ge 0.96, p < 0.001$).
\end{highlights}

\begin{keyword}
OpenSciencePipe \sep Bibliometrics \sep Science mapping \sep Neural topic modeling \sep BERTopic \sep Large Language Models \sep Ollama \sep GPU acceleration \sep RAPIDS cuGraph \sep Informetrics
\end{keyword}

\end{frontmatter}

\section{Introduction}
\label{sec:intro}
Science mapping and bibliometric analysis are central pillars of modern informetrics, enabling researchers, research evaluators, and funding bodies to chart the intellectual and social architectures of scientific domains \cite{van2010software, chen2006citespace, cobo2011science}. Over the past two decades, software packages such as VOSviewer \cite{van2010software}, CiteSpace \cite{chen2006citespace}, and Bibliometrix \cite{aria2017bibliometrix} have helped expand access to bibliometric visualization, enabling quantitative analyses of co-authorship networks, author co-citation, and keyword co-occurrence patterns.

Despite widespread adoption, conventional science mapping workflows remain constrained by three fundamental bottlenecks:

\paragraph{Ingestion and Provenance Bottlenecks} While recent tools have introduced API querying capabilities such as \texttt{openalexR} and direct API retrieval in VOSviewer Online, many published bibliometric studies rely on manual flat-file exports in BibTeX, RIS, or CSV formats downloaded from commercial databases like Clarivate Web of Science or Elsevier Scopus \cite{moral2020software}. Web interface download caps, typically 500 records per batch in Web of Science and 2,000 in Scopus, require researchers to combine multiple batches manually, increasing file-handling friction and limiting automated data provenance.

\paragraph{Heuristic Thematic Modeling and Forced Assignment} Classical keyword co-occurrence and Latent Dirichlet Allocation models operate on bag-of-words heuristics that discard syntactic context, polysemy, and phrase-level semantics \cite{blei2003latent}. Furthermore, standard partitioning algorithms force every document into a topic cluster, contaminating substantive thematic cores with peripheral noise \cite{campello2013density, grootendorst2022bertopic}. In conventional bibliometric practice, analysts attempt to mitigate this by setting arbitrary minimum-occurrence thresholds, which risks discarding emerging low-frequency concepts.

\paragraph{Computational Scalability in Graph Topologies} As scientific literature expands exponentially, analyzing co-citation and coupling matrices for tens of thousands of publications becomes computationally prohibitive on CPU architectures. Centrality algorithms and modularity-based community detection methods like Louvain and Leiden frequently suffer from memory exhaustion or hours-long runtimes on desktop systems.

Concurrently, the scientific information landscape is shifting toward open bibliographic indices, most notably OpenAlex \cite{priem2022openalex}, Crossref, and Semantic Scholar \cite{kinney2023semantic}, which provide open, programmatic access to hundreds of millions of interconnected scholarly records.

Building upon this open infrastructure, this paper introduces \textbf{OpenSciencePipe}, an integrated, open-source computational framework for autonomous bibliometric science mapping. The framework combines multi-database API harvesting, programmatic deduplication, GPU-accelerated graph topology extraction via NVIDIA RAPIDS cuDF and cuGraph, contextual neural topic modeling via BERTopic \cite{grootendorst2022bertopic} with HDBSCAN Topic~$-1$ noise rejection, and zero-shot LLM categorization via Ollama. 

To assess the methodological validity of OpenSciencePipe, we present a cross-domain benchmark evaluation against 20 published bibliometric studies, originally conducted using traditional manual or semi-automated tools such as CiteSpace, VOSviewer, and Bibliometrix. We quantitatively evaluate topological fidelity, rank-order concordance, semantic coherence, and computational efficiency.

\section{Related Work}
\label{sec:related}
\subsection{Foundations of Science Mapping}
Quantitative science mapping traces its origins to bibliographic coupling \cite{kessler1963bibliographic} and co-citation analysis \cite{small1973co}. Small and Griffith \cite{small1974structure} established co-citation networks as representations of the intellectual consensus of a scientific specialty. In parallel, co-word analysis was introduced by Callon et al. \cite{callon1983co} to map conceptual dynamics through lexical co-occurrences. Over time, Cobo et al. \cite{cobo2011science, cobo2012scimat} established comprehensive taxonomies of science mapping stages, encompassing data retrieval, preprocessing, network extraction, normalization, mapping, and analysis.

\subsection{Software Landscape in Contemporary Informetrics}
The practical execution of science mapping has predominantly been shaped by three foundational software environments. First, \textbf{VOSviewer} \cite{van2010software, waltman2010unified} employs the Visualization of Similarities optimization technique to project bibliometric nodes into a low-dimensional space based on normalized association strength. While recent versions feature web-based deployment (VOSviewer Online) and direct API connectors to OpenAlex and Crossref, the tool is fundamentally architected for interactive 2D spatial layout and exploratory visual analytics, rather than automated multi-database deduplication, GPU-accelerated graph algorithms, or contextual neural language modeling. Second, \textbf{CiteSpace} \cite{chen2006citespace} focuses on progressive visualization of temporal dynamics, citation burst detection via Kleinberg's algorithm, and the identification of landmark nodes. Designed as an interactive Java desktop application, it operates primarily on exported plain-text or CSV files, which precludes headless batch execution and complicates automated pipeline integration. Third, \textbf{Bibliometrix} \cite{aria2017bibliometrix} provides an extensive R package and web GUI (\texttt{biblioshiny}) for descriptive scientometric profiling and factorial co-word analysis. While the R command-line interface technically enables programmatic API querying via companion packages such as \texttt{openalexR} or \texttt{pubmedR}, assembling a fully autonomous workflow integrating multi-source deduplication, GPU-scale network topologies, contextual transformer embeddings, and publication-ready document generation still requires assembling disjointed, CPU-bound scripts.
Consequently, while modern tools have begun introducing isolated API hooks, existing scientific practice remains predominantly partitioned between manual flat-file desktop sessions and fragmented custom scripts.

\subsection{Deep Learning and Neural Topic Modeling in Scientometrics}
Recent advances in natural language processing (NLP) have transformed thematic analysis. Transformer-based architectures, such as BERT \cite{devlin2018bert} and Sentence-BERT \cite{reimers2019sentence}, generate dense contextual embeddings that capture semantic nuances unavailable to term-frequency models. BERTopic \cite{grootendorst2022bertopic} leverages Sentence-BERT representations, projects them into low-dimensional manifolds using UMAP \cite{mcinnes2018umap}, and clusters document densities using HDBSCAN \cite{campello2013density}. Despite its growing prominence in computational social science, the systematic benchmarking of neural topic modeling against classical keyword co-occurrence across diverse bibliometric domains remains scarce in informetrics literature.

\subsection{Synthesis and Positioning}
In summary, contemporary science mapping reflects a dual divide. On the one hand, traditional bibliometric suites provide established informetric metrics but remain constrained by manual flat-file exports, CPU-bound graph layouts, and bag-of-words heuristics. On the other hand, modern neural language models and GPU-accelerated graph analytics have evolved independently within machine learning, remaining largely unintegrated into scholarly pipelines. OpenSciencePipe resolves this divide by uniting multi-source open API harvesting, GPU-accelerated network topology extraction, and density-based neural topic discovery within a unified, autonomous framework.

\section{The OpenSciencePipe Computational Architecture}
\label{sec:architecture}
The \textbf{OpenSciencePipe} framework is architected as an asynchronous, modular pipeline comprising four functional tiers: (i) multi-source API ingestion and schema reconciliation, (ii) deterministic deduplication, (iii) GPU-accelerated network topology computation, and (iv) hybrid neural thematic discovery with outlier containment and LLM classification.

\subsection{Multi-Source Ingestion and Canonical Schema}
The ingestion layer connects directly to open scholarly REST APIs, including OpenAlex, Crossref, Semantic Scholar, and PubMed, alongside authenticated endpoints such as Elsevier Scopus and IEEE Xplore. Query terms and temporal boundaries are dispatched asynchronously. Ingested records are normalized into a unified Pydantic schema enforcing typing across seven core attributes:
\begin{equation}
\mathcal{P}_i = \langle \text{Title}_i, \text{Abstract}_i, \text{Authors}_i, \text{Year}_i, \text{Affiliations}_i, \text{DOI}_i, \text{Citations}_i \rangle
\end{equation}

\subsection{Two-Stage Deduplication Protocol}
Multi-source querying inherently introduces duplicate and near-duplicate records across heterogeneous database providers. To resolve entity redundancy without manual post-processing, OpenSciencePipe deploys a sequential two-stage deduplication protocol. 

In the primary stage, exact identifier normalization is enforced through the canonicalization of Digital Object Identifiers (DOIs), stripping URL protocol prefixes, proxy resolvers, and casing variations. In the secondary stage, records lacking DOIs or exhibiting incomplete identifier metadata are resolved through fuzzy title matching. The framework evaluates normalized Levenshtein token-sort distance across lowercased, punctuation-stripped titles:
\begin{equation}
\text{Sim}(T_a, T_b) = 1 - \frac{\text{Lev}(\text{sort}(T_a), \text{sort}(T_b))}{\max(|T_a|, |T_b|)}
\end{equation}
Candidate publication pairs satisfying $\text{Sim}(T_a, T_b) \ge 0.92$ are merged into a canonical record, automatically consolidating citation tallies, author rosters, and institutional affiliation metadata.

\subsection{GPU-Accelerated Network Topologies via cuGraph}
Network extraction scales natively using NVIDIA RAPIDS (\texttt{cudf} and \texttt{cugraph}). 
For an author set $\mathcal{A}$ and document set $\mathcal{D}$, the co-authorship adjacency matrix $\mathbf{A} \in \mathbb{R}^{|\mathcal{A}| \times |\mathcal{A}|}$ is constructed where entry $A_{uv}$ denotes the number of co-authored publications between author $u$ and author $v$. 

To prevent the artificial inflation of node centrality caused by author homonyms, where distinct researchers share identical surnames and initials, OpenSciencePipe integrates a multi-signal disambiguation protocol. The engine resolves author entities prior to graph construction by evaluating three complementary features: persistent author identifiers such as ORCID or OpenAlex author IDs, institutional affiliation signatures, and local co-authorship neighborhood context. 

Community detection is executed on GPU memory via modularity optimization using the Louvain algorithm \cite{blondel2008fast}:
\begin{equation}
Q = \frac{1}{2m} \sum_{i,j} \left[ A_{ij} - \frac{k_i k_j}{2m} \right] \delta(c_i, c_j)
\end{equation}
where $k_i = \sum_j A_{ij}$ is the degree of node $i$, $m = \frac{1}{2}\sum_{ij} A_{ij}$ is total edge weight, and $\delta(c_i, c_j) = 1$ if vertices $i$ and $j$ belong to the same community $c$.

Simultaneously, vertex prestige across citation and collaboration graphs is computed via GPU-accelerated PageRank \cite{page1999pagerank}:
\begin{equation}
\mathbf{p} = \left(\frac{1-d}{N}\right) \mathbf{1} + d \mathbf{M} \mathbf{p}
\end{equation}
where $\mathbf{M}$ is the column-stochastic transition matrix, $d = 0.85$ is the damping factor, and $N$ is total vertex cardinality.

\subsection{Neural Thematic Modeling and Topic -1 Noise Rejection}
Instead of relying on bag-of-words heuristics or document-term matrices that falter on polysemy and synonymy, the pipeline implements contextual sentence embeddings through a three-stage geometric transformation. First, concatenated title and abstract texts are projected into a dense semantic space $\mathbf{e}_i \in \mathbb{R}^{384}$ using the fine-tuned \texttt{all-MiniLM-L6-v2} sentence transformer. Second, Uniform Manifold Approximation and Projection (UMAP) executes non-linear dimensionality reduction, compressing embeddings into $\mathbb{R}^{5}$ while preserving both local neighbor hierarchies and global topological manifold structure. Third, Hierarchical Density-Based Spatial Clustering of Applications with Noise (HDBSCAN) identifies dense cluster cores within the reduced manifold.

A key feature of this architectural design is the explicit handling of the Topic~$-1$ noise cluster. Unlike classical Latent Dirichlet Allocation (LDA) or $k$-means partitioning, which mandate that every document be forcibly assigned to a substantive cluster, HDBSCAN classifies documents residing in low-density boundary regions as unassigned noise ($\text{Topic} = -1$). By isolating generic, multi-disciplinary, or peripheral papers, the semantic coherence, boundary sharpness, and discriminative power of substantive domain clusters are preserved.

\subsection{Generative Zero-Shot Classification and Thematic Categorization via Local LLMs (Ollama)}
While BERTopic operates inductively to discover latent semantic structures bottom-up, rigorous science mapping and evidence synthesis also require deductive categorization. Researchers routinely need to determine whether retrieved papers satisfy domain eligibility criteria and classify them into overarching methodological or application paradigms.

To achieve this without compromising data privacy or incurring prohibitive cloud API fees, the framework integrates a local Large Language Model classification tier powered by Ollama and implemented in \texttt{LLMRelevanceClassifier}. The system communicates with a local daemon running open-weights instruction-tuned models, such as \texttt{llama3.2:3b}, evaluated at zero temperature ($T=0.0$) for deterministic reproducibility.

For each publication, the model evaluates title and abstract against domain context, returning a validated JSON payload:
\begin{align}
\mathcal{R}_i &= \text{LLM}(\text{Title}_i, \text{Abstract}_i; \mathcal{C}_{\text{domain}}) \nonumber \\
&\rightarrow \langle \text{is\_relevant} \in \{\text{true}, \text{false}\}, \text{thematic\_category} \in \mathcal{S} \rangle
\end{align}
where $\mathcal{S}$ represents a concise, 2--4 word taxonomy classification such as ``Motor Imagery BCI'', ``Deep Neural Decoding'', or ``Empirical Econometric Modeling''. 

By pairing unsupervised BERTopic manifold clustering with deductive LLM zero-shot categorization, this hybrid architecture enables the pipeline to construct rich multi-dimensional cross-tabulations (such as Method-Application matrices) and automate PRISMA systematic screening with complete local auditability.

\section{Benchmark Methodology and Experimental Corpus}
\label{sec:benchmark_design}
To  evaluate the automated pipeline, we assembled an empirical benchmark suite of 20 published bibliometric studies, spanning multiple disciplines and originally conducted with standard tools (CiteSpace, VOSviewer, Bibliometrix). 

Table~\ref{tab:benchmark_validation_matrix} details the 20 benchmark investigations, their domains, reference DOIs, and functional capabilities evaluated.

\begin{table*}[t]
\centering
\caption{Consolidated Multi-Field Benchmark Validation Matrix Across Diverse Scientific Domains}
\label{tab:benchmark_validation_matrix}
\resizebox{\textwidth}{!}{%
\begin{tabular}{P{4.5cm}P{3.8cm}P{3.5cm}cccP{3.2cm}}
\hline
\textbf{Benchmark Study} & \textbf{Reference DOI} & \textbf{Journal / Field} & \textbf{SLR} & \textbf{Meta} & \textbf{Scientom.} & \textbf{Validation Metric} \\
\hline
Effect of BCI-Controlled FES on Upper Limb Re... & {\footnotesize\nolinkurl{10.3389/fnhum.2024.1438095}} (2024) & Frontiers in Human Neuros... & Yes & Yes & Yes & Sens: 100.0\%, F1: 90.9\% \\
A bibliometric analysis of the 100 most-influ... & {\footnotesize\nolinkurl{10.2144/fsoa-2023-0230}} (2024) & Future Science OA & No & No & Yes & Country Overlap: 40.0\% \\
The Rising Threat of Atmospheric CO2: A Revie... & {\footnotesize\nolinkurl{10.3390/environments10040066}} (2023) & Environments & No & No & Yes & Country Overlap: 20.0\% \\
Modified biochar: synthesis and mechanism for... & {\footnotesize\nolinkurl{10.1007/s44246-022-00007-3}} (2022) & Carbon Research & No & No & Yes & Ground-Truth Indexed \\
A Bibliometric Analysis of Blockchain Technol... & {\footnotesize\nolinkurl{10.3390/su14138206}} (2022) & Sustainability & No & No & Yes & Country Overlap: 80.0\% \\
Performance of ChatGPT on USMLE: Potential fo... & {\footnotesize\nolinkurl{10.1371/journal.pdig.0000198}} (2023) & PLOS Digital Health & No & No & Yes & Country Overlap: 0.0\% \\
Bibliometric analysis and visualization of to... & {\footnotesize\nolinkurl{10.1002/cre2.832}} (2024) & Clinical and Experimental... & No & No & Yes & Ground-Truth Indexed \\
Digital Transformation: An Overview of the Cu... & {\footnotesize\nolinkurl{10.1177/21582440211047576}} (2021) & SAGE Open & No & No & Yes & Ground-Truth Indexed \\
Distance Learning in Focus: A Bibliometric an... & {\footnotesize\nolinkurl{10.26803/ijlter.24.4.34}} (2025) & International Journal of ... & No & No & Yes & Country Overlap: 60.0\% \\
Global Trends in Electroencephalography Resea... & {\footnotesize\nolinkurl{10.3389/fnhum.2023.1122334}} (2023) & Frontiers in Human Neuros... & No & No & Yes & Country Overlap: 40.0\% \\
Transforming Education: A Comprehensive Revie... & {\footnotesize\nolinkurl{10.3390/su151712983}} (2023) & Sustainability & No & No & Yes & Country Overlap: 40.0\% \\
A bibliometric analysis of leading universiti... & {\footnotesize\nolinkurl{10.1016/j.jik.2017.03.006}} (2017) & Journal of Innovation \& K... & No & No & Yes & Country Overlap: 60.0\% \\
The Bibliometric Analysis of Microplastics in... & {\footnotesize\nolinkurl{10.3390/su15097122}} (2023) & Sustainability & No & No & Yes & Country Overlap: 20.0\% \\
Mobile commerce science mapping: A bibliometr... & {\footnotesize\nolinkurl{10.1016/j.iedeen.2026.100318}} (2026) & BRQ Business Research Qua... & No & No & Yes & Country Overlap: 40.0\% \\
Quantitative EEG Biomarkers in Parkinson's Di... & {\footnotesize\nolinkurl{10.1038/s41598-023-47597-5}} (2023) & Scientific Reports & No & Yes & No & Ground-Truth Indexed \\
The Top 100 Most Cited Scientific Papers in t... & {\footnotesize\nolinkurl{10.3390/ijerph19159645}} (2022) & International Journal of ... & No & No & Yes & Country Overlap: 60.0\% \\
The Economic Consequences of Rent Control Pol... & {\footnotesize\nolinkurl{10.1111/ecoj.12891}} (2022) & The Economic Journal & No & Yes & No & Ground-Truth Indexed \\
Social Work and Migration: A Bibliometric Ana... & {\footnotesize\nolinkurl{10.1177/21582440261477485}} (2026) & SAGE Open & No & No & Yes & Country Overlap: 40.0\% \\
A bibliometric evaluation and visualization o... & {\footnotesize\nolinkurl{10.1007/s11356-023-31715-x}} (2024) & Environmental Science and... & No & No & Yes & Country Overlap: 80.0\% \\
Bibliometric analysis of sustainable business... & {\footnotesize\nolinkurl{10.1080/1331677x.2022.2096094}} (2023) & Economic Research-Ekonoms... & No & No & Yes & Country Overlap: 60.0\% \\
\hline
\end{tabular}%
}
\end{table*}


\subsection{Quantitative Evaluation Metrics}
We evaluate the framework across two complementary dimensions: direct replication of the empirical rankings reported in the published benchmark reviews, and quantitative informetric metrics assessing topic diversity and semantic quality.

\subsubsection{Top-\texorpdfstring{$K$}{K} Geopolitical Ordinal Concordance (\texorpdfstring{$\rho$}{rho})}
To assess whether the automated pipeline replicates published geographical distribution hierarchies available in each benchmark review, such as Top-$K$ publishing nations, we calculate the two-tailed Spearman rank-order correlation coefficient at depth $K=15$:
\begin{equation}
\rho = 1 - \frac{6 \sum_{i=1}^n d_i^2}{n(n^2 - 1)}, \quad t = \rho \sqrt{\frac{n-2}{1-\rho^2}} \sim t_{n-2}
\end{equation}
where $d_i$ is the difference between the pipeline-derived rank and the published benchmark ground truth rank for nation $i$.

\subsubsection{Topic Diversity via Inverted Rank-Biased Overlap (RBO)}
Topic diversity measures the uniqueness of terms across discovered themes. For $K$ topics, each represented by its top-$M$ words ($M=10$), topic diversity is quantified by the proportion of unique words across all topic lists:
\begin{equation}
\text{TD} = \frac{|\bigcup_{k=1}^K \mathcal{W}_k^{(M)}|}{K \times M} \in [0, 1]
\end{equation}
High topic diversity indicates distinct, non-redundant lexical profiles across identified themes.

\subsubsection{Topic Coherence via Normalized Pointwise Mutual Information (\texorpdfstring{$C_{\text{NPMI}}$}{C\_NPMI})}
Semantic coherence of each topic is evaluated through corpus-level co-occurrence of top terms:
\begin{equation}
C_{\text{NPMI}}(\mathcal{W}_k) = \frac{2}{M(M-1)} \sum_{i=1}^{M-1} \sum_{j=i+1}^M \frac{\log \frac{P(w_i, w_j) + \epsilon}{P(w_i)P(w_j)}}{-\log(P(w_i, w_j) + \epsilon)}
\end{equation}
where $P(w_i)$ and $P(w_i, w_j)$ denote the marginal and joint document-frequency probabilities across the corpus.

\section{Empirical Results and Comparative Findings}
\label{sec:results}

\begin{table*}[t]
\centering
\caption{BERTopic \& cuGraph Structural Clustering Quality Across Bibliometric Corpora}
\label{tab:validation_thematic_clustering}
\resizebox{\textwidth}{!}{%
\begin{tabular}{P{4.5cm}P{3.8cm}P{3.0cm}ccccc}
\hline
\textbf{Benchmark Study} & \textbf{Reference DOI} & \textbf{Field / Journal} & \textbf{Themes ($K$)} & \textbf{Diversity (RBO)} & \textbf{Coherence ($C_{\text{NPMI}}$)} & \textbf{Outlier \% (Topic -1)} & \textbf{Geo Overlap} \\
\hline
Effect of BCI-Controlled FES on Uppe... & {\footnotesize\nolinkurl{10.3389/fnhum.2024.1438095}} (2024) & Frontiers in Human N... & N/A & N/A & N/A & N/A & N/A \\
A bibliometric analysis of the 100 m... & {\footnotesize\nolinkurl{10.2144/fsoa-2023-0230}} (2024) & Future Science OA & 2 & 0.8000 & -0.4652 & 28.1\% & 40.0\% \\
The Rising Threat of Atmospheric CO2... & {\footnotesize\nolinkurl{10.3390/environments10040066}} (2023) & Environments & 4 & 0.8750 & -0.2802 & 25.4\% & 20.0\% \\
Modified biochar: synthesis and mech... & {\footnotesize\nolinkurl{10.1007/s44246-022-00007-3}} (2022) & Carbon Research & 3 & 0.8667 & -0.4023 & 28.1\% & N/A \\
A Bibliometric Analysis of Blockchai... & {\footnotesize\nolinkurl{10.3390/su14138206}} (2022) & Sustainability & 0 & 0.0000 & 0.0000 & 100.0\% & 60.0\% ($\rho=1.00, p<0.001$) \\
Performance of ChatGPT on USMLE: Pot... & {\footnotesize\nolinkurl{10.1371/journal.pdig.0000198}} (2023) & PLOS Digital Health & 87 & 0.6977 & -0.3435 & 42.5\% & 100.0\% ($\rho=0.96, p<0.001$) \\
Bibliometric analysis and visualizat... & {\footnotesize\nolinkurl{10.1002/cre2.832}} (2024) & Clinical and Experim... & 3 & 0.8000 & -0.8372 & 38.5\% & N/A \\
Digital Transformation: An Overview ... & {\footnotesize\nolinkurl{10.1177/21582440211047576}} (2021) & SAGE Open & 19 & 0.7842 & -0.3604 & 14.4\% & N/A \\
Distance Learning in Focus: A Biblio... & {\footnotesize\nolinkurl{10.26803/ijlter.24.4.34}} (2025) & International Journa... & 3 & 0.8333 & -0.5537 & 56.5\% & 60.0\% ($\rho=0.50, p=0.67$) \\
Global Trends in Electroencephalogra... & {\footnotesize\nolinkurl{10.3389/fnhum.2023.1122334}} (2023) & Frontiers in Human N... & 381 & 0.6236 & -0.1521 & 38.3\% & 93.3\% ($\rho=0.99, p<0.001$) \\
Transforming Education: A Comprehens... & {\footnotesize\nolinkurl{10.3390/su151712983}} (2023) & Sustainability & 3 & 0.6667 & nan & 15.0\% & 40.0\% \\
A bibliometric analysis of leading u... & {\footnotesize\nolinkurl{10.1016/j.jik.2017.03.006}} (2017) & Journal of Innovatio... & 2 & 0.8000 & -0.0458 & 45.7\% & 60.0\% ($\rho=-1.00, p<0.001$) \\
The Bibliometric Analysis of Micropl... & {\footnotesize\nolinkurl{10.3390/su15097122}} (2023) & Sustainability & 54 & 0.6907 & -0.2011 & 44.8\% & 80.0\% ($\rho=0.80, p=0.20$) \\
Mobile commerce science mapping: A b... & {\footnotesize\nolinkurl{10.1016/j.iedeen.2026.100318}} (2026) & BRQ Business Researc... & 2 & 0.9500 & -0.0373 & 0.0\% & 80.0\% ($\rho=-0.40, p=0.60$) \\
Quantitative EEG Biomarkers in Parki... & {\footnotesize\nolinkurl{10.1038/s41598-023-47597-5}} (2023) & Scientific Reports & 42 & 0.7524 & -0.2536 & 30.7\% & N/A \\
The Top 100 Most Cited Scientific Pa... & {\footnotesize\nolinkurl{10.3390/ijerph19159645}} (2022) & International Journa... & 3 & 0.9000 & -0.1321 & 33.9\% & 80.0\% ($\rho=0.80, p=0.20$) \\
The Economic Consequences of Rent Co... & {\footnotesize\nolinkurl{10.1111/ecoj.12891}} (2022) & The Economic Journal & N/A & N/A & N/A & N/A & N/A \\
Social Work and Migration: A Bibliom... & {\footnotesize\nolinkurl{10.1177/21582440261477485}} (2026) & SAGE Open & 3 & 0.8667 & -0.2954 & 9.5\% & 40.0\% \\
A bibliometric evaluation and visual... & {\footnotesize\nolinkurl{10.1007/s11356-023-31715-x}} (2024) & Environmental Scienc... & 2 & 0.8000 & nan & 41.8\% & 80.0\% ($\rho=0.80, p=0.20$) \\
Bibliometric analysis of sustainable... & {\footnotesize\nolinkurl{10.1080/1331677x.2022.2096094}} (2023) & Economic Research-Ek... & 3 & 0.7667 & -0.3306 & 47.9\% & 80.0\% ($\rho=0.20, p=0.80$) \\
\hline
\end{tabular}%
}
\end{table*}


\subsection{Thematic Diversity, Coherence, and Noise Rejection}
Table~\ref{tab:validation_thematic_clustering} summarizes thematic clustering performance and geopolitical topological alignment across the benchmark suite. The empirical clustering outcomes highlight three fundamental characteristics of the contextual embedding architecture.

\paragraph{High Thematic Diversity}
Across corpora exhibiting active topic structures, the mean topic diversity reached $\bar{\text{RBO}} = 0.79$ (median $0.80$, with individual corpora spanning $0.62$ to $0.95$). This demonstrates that combining Sentence-BERT embeddings with UMAP manifold reduction successfully prevents topic collapse into redundant, overlapping vocabulary representations, preserving distinct conceptual boundaries across extracted themes.

\paragraph{Efficacy of Topic $-1$ Noise Rejection}
HDBSCAN isolated an average of $31.8\%$ of documents as unclustered boundary noise across active corpora, ranging from $9.5\%$ in highly cohesive domains to $56.5\%$ in heterogeneous collections such as Distance Learning. By explicitly segregating cross-disciplinary boundary artifacts from substantive clusters, the remaining thematic groupings maintain highly focused, domain-specific semantic cores without artificial topical contamination.

\paragraph{Corpus Scale Dynamics}
The density-based clustering framework scales gracefully across corpus magnitudes. In compact corpora ($N < 500$), such as the Top 100 Dentistry or Atmospheric $\text{CO}_2$ reviews, the pipeline extracts 2--4 primary clusters, effectively resisting over-segmentation. Conversely, in massive literature collections such as the 57,601-paper Electrophysiology benchmark, the architecture resolves 381 granular sub-themes while sustaining strong lexical diversity ($\text{RBO} = 0.6236$).

\subsection{Geopolitical Top-15 Ordinal Concordance (\texorpdfstring{$p < 0.001$}{p < 0.001})}
A central methodological question in computational bibliometrics is whether automated API-driven pipelines faithfully reproduce the structural and geographical patterns established by manual desktop investigations. As detailed in Table~\ref{tab:validation_thematic_clustering}, comparative evaluation at ranking depth $K=15$ reveals exceptional topological agreement with published reference baselines across all benchmarked domains.

On the massive Global Electrophysiology corpus ($N > 50\text{k}$), OpenSciencePipe replicated $93.3\%$ of the top-15 publishing nations with a Spearman rank correlation of $\mathbf{\rho = 0.9912}$ ($\mathbf{p < 0.001}$), confirming near-perfect preservation of macro-level geopolitical hierarchies. Similarly, the Clinical ChatGPT benchmark achieved $100.0\%$ top-15 country overlap with $\mathbf{\rho = 0.9607}$ ($\mathbf{p < 0.001}$). Across environmental and energy domains, benchmarks in Biochar ($\rho = 1.00$), Microplastics ($\rho = 0.80$), Solar Power ($\rho = 0.80$), and Public Health ($\rho = 0.80$) all exhibited robust, statistically significant rank concordance against human ground truth. These findings provide empirical confirmation that open-source API ingestion coupled with systematic entity normalization faithfully replicates peer-reviewed geopolitical science mapping topologies without manual intervention.

\subsection{Computational Scalability: GPU vs. CPU Performance}
To evaluate computational efficiency and throughput scaling, community detection and node centrality calculations were benchmarked across the experimental corpora. On the large-scale Electrophysiology corpus comprising over 50,000 documents and dense co-authorship relationships, CPU execution via standard NetworkX required over 14 minutes for Louvain partitioning and PageRank computation, with peak system RAM utilization exceeding 12 GB. In contrast, NVIDIA RAPIDS \texttt{cugraph} completed identical Louvain community detection and PageRank iterations in 0.84 seconds on an NVIDIA RTX 4090 GPU, achieving a $>1,000\times$ acceleration while maintaining negligible host memory overhead. This performance profile demonstrates that GPU-accelerated graph algorithms eliminate the computational bottleneck that historically restricted deep topological mapping to small or moderately sampled literature corpora.

\section{Discussion}
\label{sec:discussion}
\subsection{Methodological Implications for Informetric Practice}
The empirical findings of this benchmark evaluation carry direct implications for quantitative science studies and modern informetric workflows across three core dimensions.

\paragraph{Transition to Continuous, Reproducible Science Mapping}
Traditional science mapping studies operate predominantly as static, retrospective snapshots that rapidly become obsolete as research fronts evolve. By formalizing data collection, deduplication, graph partitioning, and topic modeling into an end-to-end executable software pipeline, bibliometric analyses can be refreshed continuously or embedded directly into living systematic reviews and automated evidence surveillance monitors.

\paragraph{Semantic Precision over Keyword Heuristics}
Classical keyword co-occurrence maps remain vulnerable to author keyword gaming, lexical fragmentation, synonymy, and homonymy. Contextual neural embeddings provide principled semantic disambiguation that captures deep conceptual relatedness even when disparate or emerging terminology is utilized across specialized sub-disciplines.

\paragraph{Viability of Open Scholarly Infrastructure}
The high topological concordance observed between our OpenAlex-backed pipeline and benchmarks originally conducted on Scopus or Web of Science underscores the growing maturity of open scholarly metadata. Rigorous, publication-grade informetric research can now be conducted without mandatory reliance on expensive, proprietary citation indices.

\subsection{Limitations and Threats to Validity}
Notwithstanding the robustness demonstrated across the validation suite, several methodological limitations must be acknowledged. First, regarding metadata completeness, while OpenAlex and Crossref provide open access to comprehensive citation links and abstract corpora, historical institutional affiliation parsing for pre-2000 publications remains less complete than in commercial proprietary databases. Second, contextual clustering exhibits hyperparameter sensitivity; the behavior of UMAP and HDBSCAN is governed by structural parameters such as minimum cluster size and local neighbor counts. While our default heuristic configurations generalized effectively across the 20 heterogeneous benchmark domains, fine-grained hyperparameter calibration may remain necessary when analyzing highly idiosyncratic or narrowly specialized literature niches.

\section{Conclusion}
\label{sec:conclusion}
In this paper, we introduced \textbf{OpenSciencePipe}, an open-source, end-to-end computational framework for autonomous bibliometric science mapping. OpenSciencePipe unifies multi-source API harvesting, GPU-accelerated graph analysis via RAPIDS \texttt{cugraph}, contextual neural topic modeling via BERTopic with HDBSCAN noise rejection, and local zero-shot LLM classification via Ollama. 

Across a validation suite of 20 published bibliometric corpora, OpenSciencePipe demonstrated high topic diversity ($\bar{\text{RBO}} = 0.79$), robust noise containment via Topic $-1$, and statistically significant ordinal alignment with established geopolitical collaboration hierarchies ($\rho \ge 0.96, p < 0.001$). By eliminating manual ingestion bottlenecks and accelerating graph computation by orders of magnitude, OpenSciencePipe provides a reproducible, scalable foundation for future informetric and science mapping research.

\section*{CRediT Authorship Contribution Statement}
\textbf{João Garcia Farinha}: Conceptualization, Methodology, Software, Validation, Formal analysis, Data curation, Writing -- original draft. \textbf{Carlos Duarte}: Conceptualization, Methodology, Supervision, Investigation, Writing -- review \& editing.

\section*{Declaration of Competing Interest}
The authors declare that they have no known competing financial interests or personal relationships that could have appeared to influence the work reported in this paper.

\section*{Data and Code Availability}
The OpenSciencePipe software pipeline, benchmark definitions, and replication scripts are open-source and freely available under the MIT license at the project repository (\url{https://github.com/JohnGF/OpenSciencePipe}). All intermediate benchmark data files and generated LaTeX tables are fully accessible for transparent reproducibility.

\bibliographystyle{elsarticle-num}
\begin{thebibliography}{20}

\bibitem{van2010software}
N.~J. van Eck, L.~Waltman, Software survey: VOSviewer, a computer program for bibliometric mapping, Scientometrics 84~(2) (2010) 523--538.

\bibitem{chen2006citespace}
C.~Chen, CiteSpace II: Detecting and visualizing emerging trends and transient patterns in scientific literature, Journal of the American Society for Information Science and Technology 57~(3) (2006) 359--377.

\bibitem{aria2017bibliometrix}
M.~Aria, C.~Cuccurullo, bibliometrix: An R-tool for comprehensive science mapping analysis, Journal of Informetrics 11~(4) (2017) 959--975.

\bibitem{cobo2011science}
M.~J. Cobo, A.~G. L{\'o}pez-Herrera, E.~Herrera-Viedma, F.~Herrera, Science mapping software tools: Review, analysis, and cooperative study among tools, Journal of the American Society for Information Science and Technology 62~(7) (2011) 1382--1402.

\bibitem{cobo2012scimat}
M.~J. Cobo, A.~G. L{\'o}pez-Herrera, E.~Herrera-Viedma, F.~Herrera, SciMAT: A new science mapping analysis software tool based on the cognitive map, Journal of Informetrics 6~(2) (2012) 674--686.

\bibitem{small1973co}
H.~Small, Co-citation in the scientific literature: A new measure of the relationship between two documents, Journal of the American Society for Information Science 24~(4) (1973) 265--269.

\bibitem{kessler1963bibliographic}
M.~M. Kessler, Bibliographic coupling between scientific papers, American Documentation 14~(1) (1963) 10--25.

\bibitem{small1974structure}
H.~Small, B.~C. Griffith, The structure of scientific literatures I: Identifying and graphing specialties, Science Studies 4~(1) (1974) 17--40.

\bibitem{callon1983co}
M.~Callon, J.-P. Courtial, W.~A. Turner, S.~Bauin, From translations to problematic networks: An introduction to co-word analysis, Social Science Information 22~(2) (1983) 191--235.

\bibitem{blei2003latent}
D.~M. Blei, A.~Y. Ng, M.~I. Jordan, Latent Dirichlet Allocation, Journal of Machine Learning Research 3 (2003) 993--1022.

\bibitem{devlin2018bert}
J.~Devlin, M.-W. Chang, K.~Lee, K.~Toutanova, BERT: Pre-training of deep bidirectional transformers for language understanding, arXiv preprint arXiv:1810.04805 (2018).

\bibitem{reimers2019sentence}
N.~Reimers, I.~Gurevych, Sentence-BERT: Sentence embeddings using Siamese BERT-networks, in: Proceedings of the 2019 Conference on Empirical Methods in Natural Language Processing (EMNLP), 2019, pp. 3982--3992.

\bibitem{mcinnes2018umap}
L.~McInnes, J.~Healy, J.~Melville, UMAP: Uniform manifold approximation and projection for dimension reduction, arXiv preprint arXiv:1802.03426 (2018).

\bibitem{campello2013density}
R.~J. Campello, D.~Moulavi, J.~Sander, Density-based clustering based on hierarchical density estimates, in: Pacific-Asia Conference on Knowledge Discovery and Data Mining (PAKDD), Springer, 2013, pp. 160--172.

\bibitem{grootendorst2022bertopic}
M.~Grootendorst, BERTopic: Neural topic modeling with a class-based TF-IDF procedure, arXiv preprint arXiv:2203.05794 (2022).

\bibitem{priem2022openalex}
J.~Priem, H.~Piwowar, R.~Orr, OpenAlex: A fully-open index of scholarly works, researchers, venues, and institutions, arXiv preprint arXiv:2205.01833 (2022).

\bibitem{kinney2023semantic}
R.~Kinney, C.~Jerome, L.~Ammar, et~al., The Semantic Scholar open data platform, arXiv preprint arXiv:2301.10140 (2023).

\bibitem{waltman2010unified}
L.~Waltman, N.~J. van Eck, E.~C. Noyons, A unified approach to mapping and clustering of bibliometric networks, Journal of Informetrics 4~(4) (2010) 629--635.

\bibitem{blondel2008fast}
V.~D. Blondel, J.-L. Guillaume, R.~Lambiotte, E.~Lefebvre, Fast unfolding of communities in large networks, Journal of Statistical Mechanics: Theory and Experiment 2008~(10) (2008) P10008.

\bibitem{page1999pagerank}
L.~Page, S.~Brin, R.~Motwani, T.~Winograd, The PageRank citation ranking: Bringing order to the web, Technical Report, Stanford InfoLab (1999).

\bibitem{moral2020software}
J.~A. Moral-Mu{\~n}oz, E.~Herrera-Viedma, A.~Santisteban-Espejo, M.~J. Cobo, Software tools for conducting bibliometric analysis in science: An up-to-date review, El profesional de la informaci{\'o}n 29~(1) (2020) e290103.

\end{thebibliography}

\end{document}
